\UnnumberedSection{Executive Summary}

Walking an 18-hole golf course imposes significant physical strain and cognitive burden on senior golfers, who must transport heavy equipment while managing club selection and ensuring no clubs are misplaced during play[cite: 3, 9]. Existing products either transport equipment without golf metrics or log shots without offering autonomous mobility[cite: 9, 10]. To bridge this gap, the Autonomous Smart Caddy Cart (ASCC) is designed as a phone-free, all-in-one solution that seamlessly follows the golfer across dry grass, automatically stops for obstacles and remote commands, keeps an active inventory of clubs in the bag, and displays real-time GPS distances to the green on an outward-facing interface[cite: 3, 9, 10].

The system integrates a Raspberry Pi 5 core controller with an array of ultrasonic/LiDAR sensors, an RFID smart-bag slot matrix, GPS processing, and a dual-motor drive platform[cite: 9, 13, 16]. Key performance targets include maintaining a safe follow distance of $1.5\text{ m} \pm 0.3\text{ m}$ at speeds up to $5.0\text{ km/h}$ (\textbf{PF-02}, \textbf{PF-03}), reacting to ground obstacles within $200\text{ ms}$ (\textbf{PF-01}), auditing 14 club slots (\textbf{IF-03}), and sustaining a continuous operating distance of $10.0\text{ km}$ on a single charge (\textbf{PF-04})[cite: 3, 12, 13]. Development progresses through seven technical milestones between October 2026 and March 2027: (1) motorized control setup, (2) sensor hardware integration, (3) RFID club tracking, (4) obstacle avoidance and emergency stop features, (5) user interface design, (6) full system integration, and (7) final testing and validation[cite: 14].

The committed end deliverable is a single functional MVP prototype demonstrated on open grass with one golfer and a bag of tagged clubs[cite: 3, 10]. The testing methodology employs staged validation—separating stationary bag trials (club identification, display accuracy) and motion trials (following dynamics, emergency stop latency $\le 50\text{ ms}$) prior to integrated outdoor course testing[cite: 3, 13, 15]. Sourcing a used motorized golf cart alongside existing hardware keeps the estimated preliminary total project cost at \$595 CAD[cite: 3, 16]. With defined subsystem leads, established resource lead times, and early risk mitigations in place, the team is fully prepared to execute the project[cite: 13, 14, 16].
\clearpage